% ============================================================
\chapter{Revisão de Python II: listas, tuplas e funções}
\label{cap:aula2}
% ============================================================

\begin{objetivos}
\begin{itemize}
  \item criar e manipular listas: índices, fatiamento, \py{append}, \py{insert}, \py{len}, \py{sum}, \py{sorted};
  \item entender a diferença entre \emph{copiar} uma lista e dar a ela um \emph{segundo nome};
  \item percorrer listas com \py{for} e construí-las com \emph{compreensão de listas};
  \item usar tuplas (sequências imutáveis) e, brevemente, dicionários;
  \item definir funções com \py{def}, \py{return}, parâmetros com valor padrão e funções \py{lambda}.
\end{itemize}
\end{objetivos}

\begin{roteiro}
0--5 min & Revisão rápida da aula anterior (\py{for} e \py{if}).\\
5--30 min & Listas: criação, índices, fatiamento, métodos, cópia \emph{versus} referência.\\
30--40 min & Percorrendo listas; compreensão de listas.\\
40--50 min & Tuplas e dicionários.\\
50--75 min & Funções: \py{def}, \py{return}, parâmetros, escopo, \py{lambda}.\\
75--90 min & Exercícios (primos, sequências, estatística simples).\\
\end{roteiro}

% ------------------------------------------------------------
\section{Listas}
% ------------------------------------------------------------

Uma \textbf{lista} é uma sequência ordenada de valores, escrita entre colchetes e separada por vírgulas. Os
elementos podem ser de qualquer tipo, inclusive outras listas.

\begin{entrada}
notas = [9.5, 7.5, 10, 8.0]
alunos = ["Davi", "Felipe", "Henrique", "Maria"]
mista = [1, "dois", 3.0, [4, 5]]
vazia = []
print(notas)
print(len(alunos), type(alunos), len(vazia))
\end{entrada}
\begin{saida}
[9.5, 7.5, 10, 8.0]
4 <class 'list'> 0
\end{saida}

\subsection{Índices}

Cada elemento ocupa uma posição (\emph{índice}). \textbf{A contagem começa em 0}. Índices negativos contam a
partir do fim: \py{-1} é o último, \py{-2} o penúltimo etc.

\begin{figure}[ht]
\centering
\begin{tikzpicture}[font=\small, cel/.style={draw=gray!70, fill=fundocod, minimum width=1.9cm, minimum height=0.75cm, font=\ttfamily}]
  \foreach \v/\i in {"Davi"/0, "Felipe"/1, "Henrique"/2, "Maria"/3} {
     \node[cel] at (\i*1.9, 0) {\v};
     \node[text=azul, font=\ttfamily] at (\i*1.9, 0.75) {\i};
     \pgfmathtruncatemacro{\n}{\i-4}
     \node[text=vermelho, font=\ttfamily] at (\i*1.9, -0.75) {\n};
  }
  \node[anchor=east, text=azul] at (-1.2,0.75) {índice};
  \node[anchor=east, text=vermelho] at (-1.2,-0.75) {índice negativo};
  \node[anchor=east] at (-1.2,0) {\texttt{alunos}};
\end{tikzpicture}
\caption{Índices positivos e negativos da lista \texttt{alunos}.}
\label{fig:indices}
\end{figure}

\begin{entrada}
print(alunos[0], alunos[2], alunos[-1])
alunos[1] = "Laura"        # listas podem ser alteradas
print(alunos)
\end{entrada}
\begin{saida}
Davi Henrique Maria
['Davi', 'Laura', 'Henrique', 'Maria']
\end{saida}

% erro
\begin{entrada}
alunos[4]
\end{entrada}
\begin{saida}
IndexError: list index out of range
\end{saida}

Uma lista com 4 elementos tem índices de 0 a 3; pedir o índice 4 é um erro muito comum.

\subsection{Fatiamento}

A expressão \py{lista[i:j]} devolve a \emph{fatia} com os elementos de índice $i, i+1,\dots, j-1$ (como no
\py{range}, o fim não entra). Também é possível indicar um passo: \py{lista[i:j:p]}. Se \py{i} for omitido, começa
do início; se \py{j} for omitido, vai até o fim.

\begin{entrada}
v = [10, 20, 30, 40, 50, 60]
print(v[1:4])      # índices 1, 2, 3
print(v[:3])       # os 3 primeiros
print(v[3:])       # do índice 3 até o fim
print(v[-2:])      # os 2 últimos
print(v[::2])      # de 2 em 2
print(v[::-1])     # invertida
\end{entrada}
\begin{saida}
[20, 30, 40]
[10, 20, 30]
[40, 50, 60]
[50, 60]
[10, 30, 50]
[60, 50, 40, 30, 20, 10]
\end{saida}

\subsection{Operações e métodos}

\begin{center}
\begin{tabular}{@{}ll@{}}
\toprule
\textbf{Comando} & \textbf{O que faz} \\
\midrule
\py{len(L)}, \py{sum(L)}, \py{min(L)}, \py{max(L)} & tamanho, soma, menor e maior elemento\\
\py{x in L} & \py{True} se \py{x} está na lista\\
\py{L + M}, \py{3 * L} & concatena listas; repete a lista 3 vezes\\
\py{L.append(x)} & acrescenta \py{x} no \textbf{fim} da lista\\
\py{L.insert(i, x)} & insere \py{x} na posição \py{i}\\
\py{L.remove(x)} & remove a primeira ocorrência de \py{x}\\
\py{L.pop()} / \py{L.pop(i)} & remove e devolve o último elemento / o de índice \py{i}\\
\py{L.index(x)}, \py{L.count(x)} & posição da primeira ocorrência de \py{x}; quantas vezes aparece\\
\py{L.sort()} / \py{sorted(L)} & ordena a própria lista / devolve uma \emph{nova} lista ordenada\\
\py{L.reverse()} & inverte a própria lista\\
\py{L.copy()} ou \py{L[:]} & cria uma cópia da lista\\
\bottomrule
\end{tabular}
\end{center}

\begin{entrada}
notas = [9.5, 7.5, 10, 8.0]
notas.append(7.0)
notas.insert(1, 6.5)
print(notas)
print("média:", sum(notas) / len(notas))
print("maior:", max(notas), " menor:", min(notas))
print(sorted(notas))        # nova lista ordenada
print(10 in notas, 5 in notas)
print([1, 2] + [3], 3 * [0])
\end{entrada}
\begin{saida}
[9.5, 6.5, 7.5, 10, 8.0, 7.0]
média: 8.083333333333334
maior: 10  menor: 6.5
[6.5, 7.0, 7.5, 8.0, 9.5, 10]
True False
[1, 2, 3] [0, 0, 0]
\end{saida}

\subsection{Cópia ou novo nome?}

Este é um ponto que confunde muitos estudantes (e que aparece no livro da SBM, Exemplo~3.6). Ao escrever
\py{B = A}, \textbf{não} se cria uma nova lista: \py{B} passa a ser apenas outro nome para a \emph{mesma} lista.
Alterar \py{B} altera \py{A}. Para obter uma lista independente, use \py{A.copy()} ou \py{A[:]}.

\begin{figure}[ht]
\centering
\begin{tikzpicture}[font=\small,
   nome/.style={draw=azul, fill=azul!8, rounded corners=2pt, minimum width=1.1cm, minimum height=0.6cm, font=\ttfamily},
   lst/.style={draw=gray!70, fill=fundocod, minimum height=0.6cm, font=\ttfamily}]
  \node[font=\sffamily\bfseries] at (1.5,1.2) {\texttt{B = A}};
  \node[nome] (A) at (0,0.4) {A}; \node[nome] (B) at (0,-0.6) {B};
  \node[lst] (L) at (3,-0.1) {[1, 2, 3, 99]};
  \draw[-{Stealth}, thick] (A) -- (L); \draw[-{Stealth}, thick] (B) -- (L);
  \node[text=vermelho, align=center, font=\scriptsize] at (1.5,-1.5) {uma só lista, dois nomes:\\\texttt{B.append(99)} muda \texttt{A}};
  \begin{scope}[xshift=7.3cm]
  \node[font=\sffamily\bfseries] at (1.5,1.2) {\texttt{B = A.copy()}};
  \node[nome] (A2) at (0,0.4) {A}; \node[nome] (B2) at (0,-0.6) {B};
  \node[lst] (L1) at (3,0.4) {[1, 2, 3]};
  \node[lst] (L2) at (3,-0.6) {[1, 2, 3, 99]};
  \draw[-{Stealth}, thick] (A2) -- (L1); \draw[-{Stealth}, thick] (B2) -- (L2);
  \node[text=verde!70!black, align=center, font=\scriptsize] at (1.5,-1.5) {duas listas independentes};
  \end{scope}
\end{tikzpicture}
\caption{Atribuir uma lista a outro nome não a copia.}
\end{figure}

\begin{entrada}
A = [1, 2, 3]
B = A             # B é outro nome para a MESMA lista
B.append(99)
print("A =", A)

A = [1, 2, 3]
C = A.copy()      # C é uma lista nova, independente
C.append(99)
print("A =", A, "  C =", C)
\end{entrada}
\begin{saida}
A = [1, 2, 3, 99]
A = [1, 2, 3]   C = [1, 2, 3, 99]
\end{saida}

% ------------------------------------------------------------
\section{Percorrendo e construindo listas}
% ------------------------------------------------------------

O \py{for} percorre diretamente os elementos de uma lista. Se também precisarmos do índice, usamos
\py{enumerate}; para percorrer duas listas lado a lado, \py{zip}.

\begin{entrada}
alunos = ["Davi", "Felipe", "Henrique", "Maria"]
notas = [9.5, 7.5, 10, 8.0]
for nome in alunos:
    print(nome, end="; ")
print()
for i, nome in enumerate(alunos):
    print(i, nome)
for nome, nota in zip(alunos, notas):
    situacao = "aprovado" if nota >= 6 else "reprovado"
    print(f"{nome:10s} {nota:5.1f}  {situacao}")
\end{entrada}
\begin{saida}
Davi; Felipe; Henrique; Maria;
0 Davi
1 Felipe
2 Henrique
3 Maria
Davi         9.5  aprovado
Felipe       7.5  aprovado
Henrique    10.0  aprovado
Maria        8.0  aprovado
\end{saida}

\begin{observacao}
A linha \py{situacao = "aprovado" if nota >= 6 else "reprovado"} é um \py{if} ``em uma linha'' (expressão
condicional): vale o primeiro valor se a condição for verdadeira e o segundo caso contrário.
\end{observacao}

\paragraph{Construindo uma lista com \texttt{append}.} Começa-se com a lista vazia e acrescenta-se um elemento a
cada volta. Exemplo: os 10 primeiros termos da sequência de Fibonacci, $F_1=F_2=1$, $F_{n}=F_{n-1}+F_{n-2}$.

\begin{entrada}
fib = [1, 1]
for n in range(8):
    fib.append(fib[-1] + fib[-2])
print(fib)
razoes = []
for k in range(1, len(fib)):
    razoes.append(fib[k] / fib[k-1])
print(razoes[-1], (1 + 5**0.5)/2)    # converge para o número de ouro
\end{entrada}
\begin{saida}
[1, 1, 2, 3, 5, 8, 13, 21, 34, 55]
1.6176470588235294 1.618033988749895
\end{saida}

\paragraph{Compreensão de listas.} É uma forma compacta, muito parecida com a notação de conjuntos da matemática.
Compare
\[
\{k^2 : k\in\{1,\dots,10\}\},\qquad \{k\in\{1,\dots,30\} : k \text{ é múltiplo de } 7\}
\]
com

\begin{entrada}
quadrados = [k**2 for k in range(1, 11)]
multiplos7 = [k for k in range(1, 31) if k % 7 == 0]
print(quadrados)
print(multiplos7)
print(sum([1/k**2 for k in range(1, 1001)]))    # aproxima pi^2/6
\end{entrada}
\begin{saida}
[1, 4, 9, 16, 25, 36, 49, 64, 81, 100]
[7, 14, 21, 28]
1.6439345666815597
\end{saida}

\begin{sintaxe}
[expressão for variável in sequência]
[expressão for variável in sequência if condição]
\end{sintaxe}

\paragraph{Listas de listas.} Uma matriz pode ser representada como uma lista de linhas. O elemento da linha $i$ e
coluna $j$ é \py{M[i][j]} (lembrando que a contagem começa em 0). Na Aula~5 veremos uma forma muito melhor de
trabalhar com matrizes (NumPy).

\begin{entrada}
M = [[1, 2, 3],
     [4, 5, 6]]
print(M[1][2], len(M), len(M[0]))
T = [[M[i][j] for i in range(2)] for j in range(3)]   # transposta
print(T)
\end{entrada}
\begin{saida}
6 2 3
[[1, 4], [2, 5], [3, 6]]
\end{saida}

% ------------------------------------------------------------
\section{Tuplas e dicionários}
% ------------------------------------------------------------

Uma \textbf{tupla} é parecida com uma lista, mas escrita com parênteses e \textbf{imutável}: depois de criada,
não pode ser alterada. É usada para dados que ``andam juntos'', como as coordenadas de um ponto.

\begin{entrada}
P = (3, 4)
Q = (0, 1, 2)
print(P[0], P[1], len(Q), type(P))
x, y = P                  # desempacotamento
print("x =", x, " y =", y)
\end{entrada}
\begin{saida}
3 4 3 <class 'tuple'>
x = 3  y = 4
\end{saida}

% erro
\begin{entrada}
P[0] = 10
\end{entrada}
\begin{saida}
TypeError: 'tuple' object does not support item assignment
\end{saida}

Para ``alterar'' uma tupla, converte-se em lista com \py{list()} (e de volta com \py{tuple()}, se necessário).

\begin{entrada}
L = list(P)
L.append(5)
print(L, tuple(L))
\end{entrada}
\begin{saida}
[3, 4, 5] (3, 4, 5)
\end{saida}

Um \textbf{dicionário} associa \emph{chaves} a \emph{valores}, como uma tabela. É escrito com chaves \texttt{\{\}}
e pares \py{chave: valor}.

\begin{entrada}
notas = {"Davi": 9.5, "Felipe": 8.5, "Maria": 8.0}
print(notas["Felipe"])
notas["Laura"] = 7.0           # nova chave
for nome, nota in notas.items():
    print(nome, nota)
\end{entrada}
\begin{saida}
8.5
Davi 9.5
Felipe 8.5
Maria 8.0
Laura 7.0
\end{saida}

\begin{center}
\begin{tabular}{@{}llll@{}}
\toprule
& \textbf{Lista} & \textbf{Tupla} & \textbf{Dicionário}\\
\midrule
Delimitador & \texttt{[ ]} & \texttt{( )} & \texttt{\{ \}}\\
Acesso & por índice \py{L[0]} & por índice \py{T[0]} & por chave \py{D["Davi"]}\\
Pode ser alterado? & sim & não & sim\\
Uso típico & coleções que mudam & pontos, pares, retornos & tabelas, cadastros\\
\bottomrule
\end{tabular}
\end{center}

\begin{maxima}
No Maxima, listas também são escritas com colchetes, \texttt{[1, 2, 3]}, mas o primeiro elemento é
\texttt{L[1]}. Em Python é \py{L[0]}. O \texttt{makelist(k\^{}2, k, 1, 10)} do Maxima corresponde a
\py{[k**2 for k in range(1, 11)]}.
\end{maxima}

% ------------------------------------------------------------
\section{Funções}
% ------------------------------------------------------------

Uma função agrupa uma sequência de comandos sob um nome, para que possam ser reutilizados. Em matemática
escrevemos $f(x)=x^2+1$; em Python:

\begin{entrada}
def f(x):
    return x**2 + 1

print(f(2), f(-1.5), f(10))
\end{entrada}
\begin{saida}
5 3.25 101
\end{saida}

\begin{sintaxe}
def nome_da_funcao(parametro1, parametro2, ...):
    """Texto opcional explicando a função (docstring)."""
    comandos
    return resultado
\end{sintaxe}

\begin{itemize}
  \item a primeira linha começa com \py{def}, tem os parâmetros entre parênteses e termina com dois-pontos;
  \item o corpo da função é \textbf{recuado};
  \item \py{return} encerra a função e devolve o resultado para quem a chamou;
  \item a função só é executada quando é \emph{chamada}, como em \py{f(2)}.
\end{itemize}

\paragraph{\texttt{return} não é \texttt{print}.} Uma função que apenas \emph{imprime} um valor não o devolve: o
resultado não pode ser usado em outras contas.

\begin{entrada}
def dobro_imprime(x):
    print(2*x)

def dobro(x):
    return 2*x

a = dobro(5) + 1
b = dobro_imprime(5)
print("a =", a, "  b =", b)
\end{entrada}
\begin{saida}
10
a = 11   b = None
\end{saida}

(\py{None} é o valor ``nada'', devolvido por funções sem \py{return}.)

\paragraph{Vários parâmetros e vários resultados.} Uma função pode receber vários valores e devolver vários
resultados (que chegam como uma tupla). Exemplo: área e perímetro de um círculo de raio $r$ e as raízes de
$ax^2+bx+c=0$.

\begin{entrada}
import math

def circulo(r):
    """Devolve a área e o perímetro de um círculo de raio r."""
    return math.pi * r**2, 2 * math.pi * r

area, perimetro = circulo(3)
print(f"área = {area:.4f}, perímetro = {perimetro:.4f}")

def raizes(a, b, c):
    delta = b**2 - 4*a*c
    if delta < 0:
        return ()                 # tupla vazia: sem raízes reais
    r1 = (-b + math.sqrt(delta)) / (2*a)
    r2 = (-b - math.sqrt(delta)) / (2*a)
    return r1, r2

print(raizes(1, -5, 6), raizes(1, 0, 1), raizes(2, 4, 2))
\end{entrada}
\begin{saida}
área = 28.2743, perímetro = 18.8496
(3.0, 2.0) () (-1.0, -1.0)
\end{saida}

(O módulo \py{math}, que fornece $\pi$ e a raiz quadrada, é o tema da próxima aula.)

\paragraph{Parâmetros com valor padrão.} Um parâmetro pode ter um valor usado quando ele não é informado.
Os argumentos também podem ser passados pelo nome.

\begin{entrada}
def soma_pa(a1, r, n=10):
    """Soma dos n primeiros termos da PA de primeiro termo a1 e razão r."""
    return sum([a1 + k*r for k in range(n)])

print(soma_pa(1, 1))          # n = 10 (padrão): 1 + 2 + ... + 10
print(soma_pa(1, 1, 100))     # 1 + 2 + ... + 100
print(soma_pa(r=2, a1=1, n=5))
\end{entrada}
\begin{saida}
55
5050
25
\end{saida}

\paragraph{Variáveis locais.} Variáveis criadas dentro de uma função só existem dentro dela (são
\emph{locais}). Isso evita que uma função ``bagunce'' as variáveis do resto do notebook.

\begin{entrada}
def quadrado(x):
    resultado = x * x
    return resultado

print(quadrado(7))
print("resultado" in dir())    # 'resultado' não existe fora da função
\end{entrada}
\begin{saida}
49
False
\end{saida}

\paragraph{Funções que usam funções.} Funções podem chamar outras funções e podem até ser passadas como
argumento. Abaixo, \py{tabela} recebe uma função \py{g} qualquer e mostra seus valores em pontos igualmente
espaçados de $[a,b]$.

\begin{entrada}
def tabela(g, a, b, n):
    h = (b - a) / n
    for k in range(n + 1):
        x = a + k*h
        print(f"{x:6.2f} {g(x):10.4f}")

tabela(f, 0, 1, 4)
\end{entrada}
\begin{saida}
  0.00     1.0000
  0.25     1.0625
  0.50     1.2500
  0.75     1.5625
  1.00     2.0000
\end{saida}

\paragraph{Funções \texttt{lambda}.} Para funções curtas, de uma expressão só, pode-se usar \py{lambda}:
\py{g = lambda x: x**3 - x} é o mesmo que definir \py{g} com \py{def} e \py{return x**3 - x}.

\begin{entrada}
g = lambda x: x**3 - x
print(g(2))
tabela(lambda t: 1/(1 + t**2), 0, 2, 2)
\end{entrada}
\begin{saida}
6
  0.00     1.0000
  1.00     0.5000
  2.00     0.2000
\end{saida}

\paragraph{Exemplo completo: números primos.} Um inteiro $n\ge 2$ é primo se não tem divisores entre 2 e $n-1$.
Na verdade basta testar divisores $d$ com $d^2\le n$: se $n=d\cdot e$ com $d\le e$, então $d^2\le n$.

\begin{entrada}
def eh_primo(n):
    if n < 2:
        return False
    d = 2
    while d * d <= n:
        if n % d == 0:
            return False      # achou um divisor: não é primo
        d += 1
    return True

print([n for n in range(2, 50) if eh_primo(n)])
print(eh_primo(97), eh_primo(2**31 - 1))
\end{entrada}
\begin{saida}
[2, 3, 5, 7, 11, 13, 17, 19, 23, 29, 31, 37, 41, 43, 47]
True True
\end{saida}

\begin{atencao}
\begin{itemize}
  \item Esquecer os parênteses na chamada: \py{f} é a função; \py{f(2)} é o valor dela em 2.
  \item Usar a função antes de executar a célula onde ela foi definida gera \py{NameError}.
  \item Não use como nome de variável ou função nomes já existentes, como \py{sum}, \py{list}, \py{max} ou
  \py{print}: o comando original deixa de funcionar até reiniciar o kernel.
\end{itemize}
\end{atencao}

\begin{professor}
O exemplo da soma dos quadrados no livro da SBM (Exemplo 3.25) tem um erro de digitação (\texttt{i + i} em vez de
\texttt{i * i}) e usa \texttt{sum} como nome de variável, o que ``apaga'' a função \texttt{sum} do Python dentro
da função. São bons pontos para discutir com a turma. A diferença entre \texttt{print} e \texttt{return} costuma
gerar muitas dúvidas: vale insistir no exemplo \texttt{dobro} \emph{versus} \texttt{dobro\_imprime}.
\end{professor}

% ------------------------------------------------------------
\section{Exercícios}
% ------------------------------------------------------------

\begin{exercicio}[manipulando listas]{ex:2-listas}
Dada \py{L = [3, 8, -2, 7, 0, 5, -6, 1]}:
(a) mostre o terceiro elemento e o último;
(b) crie a lista dos elementos positivos;
(c) crie a lista dos quadrados dos elementos;
(d) calcule a média e mostre a lista ordenada em ordem decrescente.
\end{exercicio}

\begin{exercicio}[soma de quadrados -- livro da SBM]{ex:2-quadrados}
Escreva uma função \py{soma_quadrados(n)} que devolva $1^2+2^2+\cdots+n^2$. Use-a para verificar a fórmula
$\frac{n(n+1)(2n+1)}{6}$ para $n=1,\dots,20$.
\end{exercicio}

\begin{exercicio}[média e desvio padrão]{ex:2-estat}
Escreva funções \py{media(L)} e \py{desvio(L)} que calculem a média $\bar x$ e o desvio padrão
$\displaystyle s=\sqrt{\tfrac{1}{n}\sum (x_i-\bar x)^2}$ de uma lista. Teste com as notas \py{[9.5, 7.5, 10, 8.0, 6.5]}.
\end{exercicio}

\begin{exercicio}[primos gêmeos]{ex:2-gemeos}
Usando \py{eh_primo}, liste todos os pares de primos gêmeos $(p, p+2)$ com $p<100$.
\end{exercicio}

\begin{exercicio}[MDC e MMC]{ex:2-mdc}
Escreva uma função \py{mdc(a, b)} usando o algoritmo de Euclides (Aula 1) e uma função \py{mmc(a, b)} usando
$\mathrm{mmc}(a,b)=\frac{ab}{\mathrm{mdc}(a,b)}$. Calcule o MMC de $1, 2, \dots, 10$.
\end{exercicio}

\begin{exercicio}[extra: polinômio como lista]{ex:2-poli}
Represente o polinômio $p(x)=a_0+a_1x+\cdots+a_nx^n$ pela lista \py{[a0, a1, ..., an]}. Escreva uma função
\py{avalia(coef, x)} que calcule $p(x)$ e uma função \py{deriva(coef)} que devolva a lista de coeficientes de
$p'(x)$. Teste com $p(x)=1-3x+2x^3$.
\end{exercicio}
